\color{unidarkgreen}{\subsection{Логика отключения общие}}
\color{black}

\subsubsection{Общие сведения}

\begin{enumerate}[label=\arabic{section}.\arabic{subsection}.\arabic{subsubsection}.\arabic*,  align=left, labelsep=4pt, leftmargin=0pt, itemindent=57pt]

\item
Срабатывание логики отключения основных защит формируются при возникновении одного из
следующих событий:
\begin{itemize}
\item от сигнала \textit{«Откл. от УРОВ ВН»} при отключенном положении \textit{«ОР ОВ ВН»} или при включённом положении \textit{«ОР ОВ ВН»} от \textit{«Откл. от УРОВ ОВ ВН»};
\item Откл.ДТЗ;
\item Отключение от ГЗ Т сигн;
\item Ср. ГЗ Т откл;
\item Ср. ГЗ РПН;
\item Отключение от ГЗ ЛРТ сигн;
\item Ср. ГЗ ЛРТ откл;
\item Ср.отсеч.клап. Т;
\item Ср.предохр.клап. Т;
\item Ср.предохр.клап. ЛРТ;
\item Откл. от ТЗ Т;
\item Откл. от ТЗ ЛРТ;
\item Ср. ЗПО.
\end{itemize}

\item
Действие основных защит на отсечной клапан активируется отдельной накладкой \textbf{«Действие на
ОК»}.

\item
При ее вводе формируется команда \textbf{«Закр. ОК»}. Длительность управляющего импульса задается
уставкой \textbf{«Тимп ОК»}.

\item
Срабатывание логики отключения резервных защит формируются при возникновении одного из
следующих событий:
\begin{itemize}
\item Ср. МТЗ-2 ВН;
\item Ср. МТЗ-2 НН2(СН);
\item Ср. МТЗ-2 НН1(НН);
\item Ср. ТЗНП ВН.
\end{itemize}

\item
Логика деления от МТЗ-1 НН1(НН) выводится накладкой \textbf{«ЛД НН1(НН) от МТЗ-1 НН1(НН)»}.
Срабатывание логики деления формируются при возникновении пуска МТЗ-1 НН1(НН). Так же
срабатывание контролируется выдержкой времени \textbf{«Тср МТЗ-1 на В НН1(НН)»}.

\item
Логика деления от МТЗ-1 НН2 выводится накладкой \textbf{«ЛД НН2 от МТЗ-1 НН2(СН)»}. Срабатывание
логики деления формируются при возникновении пуска МТЗ-1 НН2. Так же срабатывание
контролируется выдержкой времени \textbf{«Тср МТЗ-1 на В НН2»}.

\item
Логика деления от ТЗНП-1 РЗН1 выводится накладкой \textbf{«ЛД НН1(НН) от ТЗНП-1 РЗН1»}.
Срабатывание логики деления формируются при возникновении пуска ТЗНП-1 РЗН1. Так же
срабатывание контролируется выдержкой времени \textbf{«Тср ТЗНП-1 РЗН1 на В НН1(НН)»}.

\item
Логика деления от ТЗНП-2 РЗН2 выводится накладкой \textbf{«ЛД НН2 от ТЗНП-1 РЗН2»}. Срабатывание
логики деления формируются при возникновении пуска ТЗНП-1 РЗН1. Так же срабатывание
контролируется выдержкой времени \textbf{«Тср ТЗНП-1 РЗН2 на В НН2»}.

\begin{figure}[H]
\centering
\includegraphics[width=0.8\textwidth,height=0.8\textheight,keepaspectratio]{\fbpath/011.9902 ЛО общ/img/LO_gen.pdf}
\captionof{figure}{Функциональная схема общей логики отключения}\label{pic:LO_gen}
\end{figure}


\small
\begin{longtable}{|>{\centering\arraybackslash}m{5.3cm}|>{\centering\arraybackslash}m{4.2cm}|>{\centering\arraybackslash}m{1cm}|>{\centering\arraybackslash}m{3.3cm}|>{\centering\arraybackslash}m{1.8cm}|}
% Первая страница: полный заголовок
\caption{Перечень уставочных параметров функции общей логики отключения\hfill\vspace{-0.5\baselineskip}}\label{desc:tbl1}\\
\hline
\rowcolor{gray!30}
Наименование уставки & Диапазон & Шаг & Значение по умолчанию & Единица измерения \\
\hline
\endfirsthead

% Промежуточные страницы: «Продолжение таблицы»
\caption*{\hspace{3pt}\emph{Продолжение таблицы \ref{desc:tbl1}}\hfill\ }\vspace{-0.5\baselineskip} \\
\hline
\rowcolor{gray!30}
Наименование уставки & Диапазон & Шаг & Значение по умолчанию & Единица измерения \\
\hline
\endhead

%===m>ДТЗ|-
%===m
\end{longtable}
\normalsize	

\end{enumerate}